\input{preamble}

\title{Section 6.6 --- Function Transformations}

\begin{document}
\maketitle

\section{Single-Step Transformations}
In each figure below, the given graph is the graph of $y=f(x)$. Graph the function in the question on the same axes. You may find it helpful to use the marked points as reference points.
\begin{smenum}
\mitemxx
{$g(x)=2f(x)$\\
\begin{GridSquareFitFilled}(2in)
\draw[thick] (-2,-1) -- (2,2) -- (4,-2);
\Cdot{-2,-1};
\Cdot{2,2};
\Cdot{4,-2};
\end{GridSquareFitFilled}}
{$g(x)=f(-x)$\\
\begin{GridSquareFitFilled}(2in)
\draw[thick] (-4,3) -- (3,-4);
\Cdot{-4,3};
\Cdot{3,-4};
\end{GridSquareFitFilled}}
\mitemxx
{$g(x)=f(x)-3$\\
\begin{GridSquareFitFilled}(2in)
\draw[thick] (-3,0) arc (180:0:3);
\Cdot{-3,0};
\Cdot{0,3};
\Cdot{3,0};
\end{GridSquareFitFilled}}
{$g(x)=-f(x)$\\
\begin{GridSquareFitFilled}(2in)
\draw[thick] (4,4) -- (0,0) -- (-3,3);
\Cdot{-3,3};
\Cdot{0,0};
\Cdot{4,4};
\end{GridSquareFitFilled}}
\mitemxx
{$g(x)=f(x)+2$\\
\begin{GridSquareFitFilled}(2in)
\draw[thick] (-3,3) parabola bend (1,-1) (3,0);
\Cdot{-3,3};
\Cdot{1,-1};
\Cdot{3,0};
\end{GridSquareFitFilled}}
{$g(x)=f(x+1)$\\
\begin{GridSquareFitFilled}(2in)
\draw[thick] (-3,-3) -- (-1,1) -- (3,1) -- (4,2);
\Cdot{-3,-3};
\Cdot{-1,1};
\Cdot{3,1};
\Cdot{4,2};
\end{GridSquareFitFilled}}
\mitemxx
{$g(x)=\frac{2}{3}f(x)$\\
\begin{GridSquareFitFilled}(2in)
\draw[thick] (-1,-3) -- (4,3);
\Cdot{-1,-3};
\Cdot{4,3};
\end{GridSquareFitFilled}}
{$g(x)=f(\frac{1}{3}x)$\\
\begin{GridSquareFitFilled}(2in)
\draw[thick] (-1,-2) parabola bend (0,3) (1,-2);
\Cdot{-1,-2};
\Cdot{0,3};
\Cdot{1,-2};
\end{GridSquareFitFilled}}
\mitemxx
{$g(x)=f(2x)$\\
\begin{GridSquareFitFilled}(2in)
\draw[thick] (-4,1) arc (-180:0:4);
\Cdot{-4,1};
\Cdot{0,-3};
\Cdot{4,1};
\end{GridSquareFitFilled}}
{$g(x)=f(x-3)$\\
\begin{GridSquareFitFilled}(2in)
\draw[thick] (-3,-2) -- (1,4);
\Cdot{-3,-2};
\Cdot{1,4};
\end{GridSquareFitFilled}}
\end{smenum}

\section{Multiple Transformations}
In each figure below, the given graph is the graph of $y=f(x)$. Graph the function in the question on the same axes. You may find it helpful to use the marked points as reference points. \emph{You will want to write out your work on a separate sheet of paper.}
\begin{smenum}
\mitemxx
{$g(x)=-2f(\frac{2}{3}x+2)+3$\\
\begin{GridSquareFitFilled}[11]
\draw[thick] (-4,-4) -- (2,3) -- (6,-2);
\Cdot{-4,-4};
\Cdot{2,3};
\Cdot{6,-2};
\end{GridSquareFitFilled}}
%\mitemx
{$g(x)=\frac{2}{3}f(-2x-5)-1$\\
\begin{GridSquareFitFilled}[11]
\draw[thick] (-9,-3) arc (180:0:6);
\Cdot{-9,-3};
\Cdot{-3,3};
\Cdot{3,-3};
\end{GridSquareFitFilled}}
\end{smenum}
\clearpage
\section{Domain and Range}
\begin{senum}
\item If $g(x)=3 f(-2x-3)+4$ and $\dom f=(-5,\infty)$, what is $\dom g$?
\item If $h(x)=\frac{1}{3} g(x+1)-7$ and $\rng g=(-\infty,3]\cup[5,\infty)$, what is $\rng h$?
\item If $k(x)=2 f\left(7-\frac{2}{5}x\right) - 6$, $\dom f=(0,1)$, and $\rng f=(-\infty,\infty)$, what are $\dom k$ and $\rng k$?
\end{senum}
\section{Symmetry}
Classify each function as even, odd, or neither.
\begin{smenum}
\maitemxxxx
{f(x)=x^2\mlbl{firstpoly}}
{f(x)=(x-7)^2}
{f(x)=x^5}
{f(x)=2x}
\maitemxxxx
{f(x)=x^6-3x^2+9}
{f(x)=x^3+x^2}
{f(x)=x^3+6}
{f(x)=x^9-x^3+8x\mlbl{lastpoly}}
\maitemxxxx
{f(x)=\sqrt{x}}
{f(x)=\sqrt[5]{x}}
{f(x)=3|2x|-5}
{f(x)=|x-5|}
\mitemx
{Based on your answers to \nnref{firstpoly}{lastpoly}, how can we decide if a polynomial is an even function, an odd function, or neither? (Try to answer this question on your own before looking at the answer key.)}
\end{smenum}
\section{Using the Library}
For each function below:
\begin{clist}
\item Identify the base function (call it $f(x)$) and the family of functions it belongs to.
\item Rewrite the function in terms of $f(x)$. (We usually won't do this explicitly in practice.)
\item Sketch a graph of the function, including any asymptotes and clearly marking any special points, using a separate sheet of graph paper or the other side of this page.
\item Give the domain and range of the function.
\end{clist}
\begin{smenum}
\maitemxx
{g(x)=-3x+5}
{g(x)=\dfrac{5}{x^2-6x+9}}
\maitemxx
{g(x)=\dfrac{3}{x-2}}
{g(x)=|x-5|+2}
\maitemxx
{g(x)=4}
{g(x)=2(x-1)^3+1}
\maitemxx
{g(x)=\log_{1/3}(2x+3)+1}
{g(x)=-\left(\frac{1}{2}\right)^{x+2}+1}
\maitemxx
{g(x)=2(x-2)+1}
{g(x)=-3(x+4)^2+2}
\maitemxx
{g(x)=\dfrac{3}{2x^2}-1}
{g(x)=3\sqrt[3]{4-x}-1}
\maitemxx
{g(x)=\dfrac{6x-7}{2x-3}}
{g(x)=2\sqrt{-\frac{1}{2}x+3}}
\maitemxx
{g(x)=2x^2-8x+5}
{g(x)=3^{\frac{1}{3}x}-5}
\maitemxx
{g(x)=\log_2(x-3)+4}
{g(x)=-\left|4-\frac{1}{3}x\right|+2}
\end{smenum}
%\section*{Messing with Scaling}
%For the functions below:
%\begin{clist}
%\item Identify the base function (call it $f(x)$) and the family of functions it belongs to.
%\item Use algebra rewrite the function as $a\cdot f\left(\frac{1}{b} x+c\right)+d$ where $a$ and $b$ are both integers. This will cause your key points to end up as integers.
%\item Sketch a graph of the function, including any asymptotes and clearly marking any special points.
%\item Give the domain and range of the function.
%\end{clist}
%\begin{smenum}
%\maitemxx
%{g(x)=3^{2x-4}-2}
%{g(x)=-\frac{3}{2}\sqrt{x+1}}
%\maitemxx
%{g(x)=\frac{1}{3} x^2}
%{g(x)=\frac{2}{3}(x-2)+5}
%\end{smenum}
\graphpage
\end{document}